\section{Discussion}\label{sec:discussion}

\subsection{What the Ablation Actually Proves}\label{subsec:what_ablation_proves}

The headline subtype number, 99.51\,\%, is the easy story. The harder
story, and the one we believe matters more for the field, is the
\textbf{role-complementarity principle}. The \texttt{bam\_triplet} cell and the
independent v2-EMA negative result, two pairings of spatial-role attention
modules reached from disjoint module sets, converge on the same 98.36\,\%
subtype ceiling, each falling below the no-attention control. It is unlikely
that two independently trained spatial-spatial configurations would land on the
same accuracy to four significant figures by chance; the probability of such a
convergence under a pure-noise null is low. We do not over-read two data points,
however: the more decisive evidence is mechanistic rather than combinatorial.
Both configurations regress \emph{beneath} the no-attention control, an
outcome a harmless redundant module could not produce, since an uninformative
module would at worst leave the representation unchanged. The convergence is
therefore better read as the signature of a systematic spatial-role conflict
than as a coincidence, and the multi-seed follow-up of
Section~\ref{subsec:future} is the proper remedy for the small sample. The
contrast with \texttt{bam\_kan} (BAM paired with the purely channel-wise KAN,
99.15\,\%) and with \texttt{triplet\_kan} (99.45\,\%) sharpens the point: a
non-spatial second module is complementary, a spatial one collides.

The consequence is that the present paper contributes not only a model but a
\emph{design rule}. When attention is added to an oral-disease classifier of
this size and on this kind of data, the modules should be paired so that they do
not both exercise the spatial role; pairing a spatial module with a purely
channel-wise module is complementary, whereas pairing two spatial modules
regresses the hub below the no-attention control. When a candidate second module
would duplicate the spatial role, the more reliable choice is to drop it. This
reframes attention-hub design from the prevailing "stack modules and report the
best configuration" practice into a measurement-driven procedure in which each
component is justified against the role taxonomy of Section~\ref{sec:prelim}.

\subsection{Comparison to Prior Multi-Class Systems}\label{subsec:vs_prior}

MODC-SET~\cite{ref22} (Section~\ref{sec:related}), the closest seven-class
result in the literature, reaches 99.32\,\% overall accuracy via a three-backbone
XGBoost ensemble. The proposed model attains 99.51\,\% subtype accuracy with a
single-stage 4.79~M parameter architecture. Because MODC-SET is evaluated on a
different seven-class collection, this is not a controlled head-to-head
comparison; we report it to situate the magnitude of the present result, not to
claim a win on a shared benchmark, and the 0.19-point delta should not be read
as a controlled accuracy advantage. The substantive contrast is methodological
rather than numerical: MODC-SET attains its accuracy by fusing features across
three large backbones, with no published ablation justifying which backbones to
ensemble or why XGBoost should serve as the meta-classifier, whereas the
accuracy of the proposed model emerges from a five-stage backbone whose only
specialized component is justified branch by branch. The two systems thus
demonstrate that high accuracy on seven-class oral disease data is attainable
by very different means; the present work shows that an efficient single-model
alternative can be reached in a principled, ablation-driven way rather than by
ensemble scale. Accordingly, the "highest subtype accuracy" claim made
throughout this paper is restricted to the controlled nine-baseline benchmark
trained here under one matched recipe; it is not asserted as a record against
the broader literature, where results obtained on different datasets and under
different protocols are not directly comparable.

Relative to the binary-screening literature~\cite{ref5,ref16,ref17}, the
proposed model maintains binary-task parity (99.06\,\% against the leading
EfficientNet~V2-B2 at 99.36\,\%) at less than half the parameter count,
while additionally delivering the seven-way subtype categorization that the
binary screeners do not address. The comparison with the segmentation
literature~\cite{ref15,ref26} is one of complementary scope rather than of
competing numbers: those systems localize lesion extent at pixel granularity,
whereas the present model identifies the disease class. The two questions are
distinct, and Section~\ref{subsec:future} returns to the prospect of combining
them.

\subsection{Why the Binary Head Plateaus}\label{subsec:binary_plateau}

Across the strongest baselines (ResNet50, Inception~V3, and
EfficientNet~V2-B2), the binary head sits within a narrow 99.1--99.4\,\%
corridor. We interpret this as a property of the data rather than of the
model: the visual gap between benign and malignant lesions is sufficiently wide
that any reasonable backbone trained from scratch resolves it almost
completely. Pushing the binary head past 99.5\,\% would require either a larger
dataset that admits harder edge cases or a task-targeted augmentation strategy;
neither is in scope here, and we do not claim a binary-task improvement.

The subtype head behaves differently. It requires the model to differentiate
among visually similar mucosal conditions, most notably MC versus OC, two
malignant carcinomas whose presentations overlap, and this is exactly where
attention-hub design has measurable leverage. The role-complementarity
principle yields its gain on the subtype head and not on the binary head
because the subtype task is the one whose residual errors are
representational rather than already saturated. The plateau on the binary head
is thus consistent with, rather than a counterexample to, the central claim of
the paper.

A clarification on the role of the dual-head construction is warranted here.
The dual-head, multi-task design is a \emph{data-utilization} mechanism: the
masked subtype loss is the means by which the binary-only DS1 collection can
contribute to the same training run as the seven-class DS2 collection, rather
than being a device claimed to raise subtype accuracy. We do not run a
single-task control (subtype-head-only training without DS1) and therefore make
no claim that joint multi-task training improves the subtype head over an
equivalent single-task model. The role-complementarity gains reported in this
paper are attributable to the AttentionHub topology, not to the multi-task
objective; isolating the contribution of joint training is left to future work,
and is recorded among the limitations below.

\subsection{Limitations}\label{subsec:limitations}

We acknowledge six limitations of the present study.

\emph{(a) Single-seed evaluation.} Each model is trained once with seed~42. For
the 0.06 percentage-point margin between AttentionHub-v2 and the best v1 cell
(\texttt{triplet\_kan}), this margin lies within the noise band of a single
run. We do not claim that v2 is \emph{statistically} better than
\texttt{triplet\_kan}; we claim that v2 is the \emph{principled} design that
the ablation forces, and that under one matched run it sets the highest subtype
number in the study.

\emph{(b) Dataset scale.} The training set is approximately 5\,500 images
across seven subtypes, small relative to natural-image benchmarks though
large relative to most oral-disease studies. We compensate through from-scratch
training (so that no ImageNet prior is smuggled in), augmentation, early
stopping, and replicated evaluation, but external validation on independent
clinics is needed before clinical deployment.

\emph{(c) Image-level only.} This work performs classification, not pixel-level
segmentation. Segmentation is a complementary direction
(Section~\ref{subsec:vs_prior}); the present model does not output lesion
masks. Bounding-box detection and pixel masks are left as future work.

\emph{(d) Training-batch timing is approximate.} The \texttt{batch\_time\_ms}
field in \texttt{performance\_metrics.json} is estimated from test-time
inference rather than measured during training. The latency numbers themselves
(mean, P50, P95) are measured correctly; only the derived "estimated epoch
time" reported alongside them is an approximation.

\emph{(e) Interpretability evidence is qualitative.} Section~\ref{sec:xai}
presents Grad-CAM++ and LIME panels and a visual alignment ranking of the
models, but the dataset carries no lesion masks and no quantitative
interpretability metric (for example deletion/insertion AUC or
pointing-game accuracy against annotated regions) is computed. The
cross-model interpretability ordering should therefore be read as indicative
rather than as a measured, falsifiable result.

\emph{(f) The dual-head contribution is not isolated.} As discussed in
Section~\ref{subsec:binary_plateau}, the dual-head multi-task design is a
data-utilization mechanism that lets the binary-only DS1 collection contribute
to training. We do not run a subtype-only single-task control, so we do not
isolate the effect of joint multi-task training on the subtype head; whether
joint training improves, leaves unchanged, or mildly degrades subtype accuracy
relative to single-task training is left to future work.

\subsection{Clinical Significance and Deployment
Considerations}\label{subsec:clinical}

The 99.51\,\% subtype accuracy reported in this study is not a clinical claim.
It is a \emph{bench-test} result on a held-out portion of a single curated
dataset, and three distinct gaps separate it from clinical deployment.

\emph{(a) Distributional gap.} The dataset's image-acquisition conditions,
patient demographics, and intra-oral-anatomy coverage are not characterized in
the published metadata. A deployment-grade screening system would require
multi-site evaluation across patient populations, imaging devices, and lighting
conditions. The proposed model's small parameter count and low inference cost
make such an evaluation tractable on commodity hardware, but the evaluation
itself has not been performed here.

\emph{(b) Decision-context gap.} The model produces a posterior over seven
classes; clinical decisions require additional inputs (history, palpation,
biopsy result) that the model does not consume. Any deployment must
therefore be framed as \emph{decision-support} rather than
\emph{decision-making}, with a clinician retaining final authority over
diagnosis and management.

\emph{(c) Calibration gap.} This work reports accuracy and F1 but does not
report calibration metrics such as Expected Calibration Error or Brier score.
For triage-style screening, where the most useful output is often a
calibrated probability that the lesion warrants biopsy, calibration becomes
more important than accuracy. We flag this as the most consequential metric the
present study does not report, and we treat its evaluation as an open
follow-up.

The proposed model is therefore released as a research artifact and as evidence
for the role-complementarity principle, not as a deployable diagnostic tool.

\subsection{Future Work}\label{subsec:future}

Five directions follow naturally from the present study.

\begin{enumerate}
  \item \textbf{Third attention slot.} Extending the role-complementarity
  principle to a \emph{third} attention slot would require a new ablation grid
  testing whether a third disjoint-role module (for example a
  frequency-domain attention, a global-context module, or a graph-based
  reasoning head) compounds the gain or whether the cascade saturates at two
  complementary modules. The principle as stated does not predict the answer.
  \item \textbf{Backbone transferability.} The v2 hub is small enough to drop
  into the larger EfficientNet~V2 family or into ResNet-style backbones;
  testing whether the same principle holds when Stage~4 sits inside a deeper or
  wider network would establish whether the rule is dataset-specific,
  architecture-specific, or general.
  \item \textbf{Multi-seed statistical evaluation.} A multi-seed run (for
  example five independent random seeds per ablation cell) would convert the
  present single-run table into a mean $\pm$ standard-deviation table, making
  within-cell variability explicit and supporting paired-sample significance
  tests between v2, \texttt{triplet\_kan}, and \texttt{full}. The principle
  predicts that the \emph{ordering} of the cells will be preserved; an
  empirical test of that prediction would strengthen the central methodological
  claim.
  \item \textbf{Hybrid classification and segmentation head.} Pairing the
  present classifier with a lightweight segmentation head that shares Stage~5
  features would recover the pixel-level information targeted by
  segmentation-only papers~\cite{ref15,ref26} while retaining the
  categorization capability targeted by classification-only
  papers~\cite{ref19,ref22}. Stage~5's 192-channel output at $7\times7$ spatial
  resolution is well suited to a small upsampling decoder.
  \item \textbf{Calibration and uncertainty.} As noted in
  Section~\ref{subsec:clinical}, calibration metrics and predictive uncertainty
  are the most consequential evaluation axes the present study does not cover.
  A follow-up applying temperature scaling, Monte-Carlo dropout, or deep
  ensembling on top of the v2 model would close this gap.
\end{enumerate}
